\section{PuzzleLemmas2}
\textit{Source: \texttt{Quadratic/Complex/PuzzleLemmas2.lean}}

This module contains the formal definitions and proofs.
\subsection*{para\_dynamical\_correspondence}
Correspondence between parameter and dynamical pieces.

\subsection*{parameter\_shrink\_ax}
The Correspondence Principle:
    If the dynamical pieces shrink to a point, the parameter pieces shrink to a point.

\subsection*{para\_puzzle\_piece\_open}
A holomorphic motion of a set E over the unit disk D. -/
structure HolomorphicMotion (E : Set ℂ) where
  /-- The motion map t ↦ z ↦ f(t, z) -/
  f : ℂ → ℂ → ℂ
  /-- At time 0, it is the identity -/
  h_zero : ∀ z ∈ E, f 0 z = z
  /-- For each fixed time in the unit disk, it is injective on E -/
  h_inj : ∀ t ∈ Metric.ball 0 1, Set.InjOn (f t) E
  /-- For each fixed z in E, it is holomorphic in time on the unit disk -/
  h_holo : ∀ z ∈ E, DifferentiableOn ℂ (fun t ↦ f t z) (Metric.ball 0 1)

/-- Slodkowski's Theorem (Generalized Lambda Lemma).
    "Every holomorphic motion f : D × E → ℂ of an arbitrary subset E of ℂ can be
    extended to a holomorphic motion F : D × ℂ → ℂ (that is F|D×E = f) of ℂ,
    parametrized by the same unit disc D."
    See: [Slodkowski, Holomorphic motions and polynomial hulls, Theorem 1.3] <https://www.ams.org/journals/proc/1991-111-02/S0002-9939-1991-1037218-8/>
    Local Reference: `refs/S0002-9939-1991-1037218-8.pdf` -/
axiom slodkowski_theorem {E : Set ℂ} (h : HolomorphicMotion E) :
    ∃ H : HolomorphicMotion Set.univ,
      ∀ t ∈ Metric.ball 0 1, ∀ z ∈ E, H.f t z = h.f t z

/-- Axiom: The boundary of a puzzle piece moves holomorphically.
    This axiom provides the existence of a holomorphic motion on the boundary of the puzzle piece,
    which serves as the input for Slodkowski's Extension Theorem.

    **Definition** [Slodkowski, p. 347]:
    "A holomorphic motion of E in C, parametrized by the unit disc D, is a map f: D × E → C
    such that (a) for any fixed w ∈ E, the map z ↦ f(z, w): D → C is holomorphic;
    (b) for any fixed z ∈ D, the map w ↦ f(z, w) is one-to-one; and (c) f₀ is the identity map on E."

    In this context:
    *   `E` is the boundary of the puzzle piece (or a neighborhood of it).
    *   The parameter `z` (or `t`) corresponds to the quadratic parameter `c`.
    *   This motion is constructed using the Böttcher coordinate, which depends holomorphically on `c`.

    Ref: `refs/S0002-9939-1991-1037218-8.pdf` (Slodkowski, Holomorphic motions and polynomial hulls). -/
axiom puzzle_boundary_motion_exists (n : ℕ) (c₀ : ℂ) (hc₀ : c₀ ∈ ParaPuzzlePiece n) :
    ∃ (r : ℝ) (_ : 0 < r) (E : Set ℂ) (h : HolomorphicMotion E),
      -- The motion is defined for parameters c in D(c₀, r) via a rescaling map
      -- ψ : D(0, 1) → D(c₀, r)
      -- And this motion preserves the "puzzle membership" property in the sense that:
      -- If H is an extension of h to the plane (guaranteed by Slodkowski),
      -- then for any t ∈ D, the corresponding parameter c_t is in ParaPuzzlePiece n.
      ∀ (H : HolomorphicMotion Set.univ),
        (∀ t ∈ Metric.ball 0 1, ∀ z ∈ E, H.f t z = h.f t z) →
        ∀ t ∈ Metric.ball 0 1, (c₀ + r * t) ∈ ParaPuzzlePiece n

/-- Parameter puzzle pieces are open sets.
    This is proved using Slodkowski's Theorem which ensures that the holomorphic motion
    of the dynamical plane implies structural stability of open sets in the parameter plane.

\subsection*{para\_puzzle\_piece\_basis}
Parameter puzzle pieces form a basis of neighborhoods if they shrink to a point.

\subsection*{para\_puzzle\_piece\_inter\_mandelbrot\_connected}
Parameter puzzle pieces intersected with the Mandelbrot set are connected.

